\documentclass[12pt,a4paper]{article}

\usepackage[T1]{fontenc}
\usepackage[utf8]{inputenc}
\usepackage[polish]{babel}
\usepackage{lmodern}
\usepackage{geometry}
\usepackage{hyperref}
\usepackage{booktabs}
\usepackage{float}

\geometry{margin=2.5cm}
\hypersetup{
    colorlinks=true,
    linkcolor=black,
    urlcolor=blue
}

\title{Sprawozdanie z projektu 4\\Własne środowisko Gymnasium}
\author{Autorzy: Marcin Wolder, Stanisław Wojtas}
\date{\today}

\begin{document}

\maketitle

\section{Cel projektu}
Celem projektu było przygotowanie własnego środowiska zgodnego z biblioteką Gymnasium, a następnie zbudowanie agenta, który potrafi skutecznie rozwiązywać zadanie występujące w tym środowisku. Jako temat wybrano uproszczoną wersję gry \textit{Crossy Road}. Zadaniem agenta jest przeprowadzenie postaci przez kolejne pasy ruchu samochodowego bez kolizji i dotarcie do wyznaczonego celu.

Projekt spełnia wymagania zadania, ponieważ obejmuje:
\begin{itemize}
    \item autorskie środowisko Gymnasium,
    \item implementację agenta uczącego się gry,
    \item eksperymenty porównujące wyuczoną politykę z zachowaniem losowym,
    \item tryb graficzny oparty na \texttt{pygame}.
\end{itemize}

\section{Opis środowiska}
Środowisko zostało zaimplementowane jako klasa \texttt{CrossyRoadEnv}. Plansza ma szerokość 10 pól i wysokość co najmniej 14 wierszy. Gracz rozpoczyna na bezpiecznym dolnym pasie, a jego celem jest osiągnięcie wiersza \texttt{goal\_distance}. Pomiędzy startem a metą znajdują się pasy ruchu oraz losowo generowane bezpieczne odcinki.

Pełna przestrzeń akcji środowiska jest dyskretna i zawiera pięć akcji:
\begin{itemize}
    \item ruch w górę,
    \item ruch w dół,
    \item ruch w lewo,
    \item ruch w prawo,
    \item oczekiwanie.
\end{itemize}

Samochody poruszają się po pasach z różnymi prędkościami i w obu kierunkach. Ich pozycje są aktualizowane w sposób ciągły, dlatego przestrzeń obserwacji nie została zbudowana jako mała siatka dyskretna, lecz jako wektor cech typu \texttt{Box}. W obserwacji znajdują się między innymi:
\begin{itemize}
    \item znormalizowana pozycja gracza,
    \item informacja, czy ruch do przodu jest aktualnie ryzykowny,
    \item informacja, czy bieżąca pozycja za chwilę stanie się niebezpieczna,
    \item dla każdego pasa ruchu cechy najbliższego samochodu: względna odległość, prędkość i kierunek,
    \item lokalna mapa zajętości wokół gracza.
\end{itemize}

Tak zdefiniowana obserwacja pozwala agentowi podejmować decyzje na podstawie stanu ruchu drogowego, a nie tylko aktualnej pozycji na planszy. Jest to istotne, ponieważ poprawna gra wymaga przewidywania nadjeżdżających samochodów i czekania na bezpieczną lukę.

\section{Funkcja nagrody}
Funkcja nagrody została zaprojektowana tak, aby promować postęp oraz zachowania strategiczne. Agent otrzymuje:
\begin{itemize}
    \item dodatnią nagrodę za przesuwanie się w stronę celu,
    \item dużą nagrodę za ukończenie poziomu,
    \item silną karę za kolizję,
    \item niewielką karę krokową za każdy ruch,
    \item dodatkowe kary za ruch wstecz, zbędne ruchy boczne i zbyt długie stanie w miejscu,
    \item premię za czekanie wtedy, gdy ruch do przodu byłby niebezpieczny,
    \item karę za wejście w ryzykowną sytuację mimo zbliżającego się zagrożenia.
\end{itemize}

Takie kształtowanie nagrody miało dwa cele. Po pierwsze, agent miał nauczyć się konsekwentnie zbliżać do mety. Po drugie, jego zachowanie miało wyglądać na celowe: wejście na jezdnię powinno następować wtedy, gdy jest to względnie bezpieczne, a nie losowo.

\section{Algorytm uczący}
Do rozwiązania problemu wykorzystano algorytm DQN z biblioteki Stable-Baselines3. Sieć była trenowana przez 100\,000 kroków. Zastosowano między innymi następujące parametry:
\begin{itemize}
    \item \texttt{learning\_rate = 3e-4},
    \item \texttt{buffer\_size = 100000},
    \item \texttt{learning\_starts = 5000},
    \item \texttt{batch\_size = 64},
    \item \texttt{gamma = 0.99},
    \item \texttt{train\_freq = 4},
    \item \texttt{target\_update\_interval = 500}.
\end{itemize}

W aktualnej wersji treningu użyto pełnej przestrzeni pięciu akcji, czyli \texttt{up}, \texttt{down}, \texttt{left}, \texttt{right} oraz \texttt{wait}. Dzięki temu agent może nie tylko przesuwać się naprzód, ale także wykonywać manewry korygujące pozycję i reagować na lokalny układ ruchu. Jest to ustawienie trudniejsze od wariantu z ograniczoną liczbą akcji, ponieważ powiększa przestrzeń decyzji i zwiększa ryzyko uczenia nieefektywnych zachowań bocznych.

Podczas treningu środowisko było okresowo ewaluowane, a najlepszy model zapisywano do pliku \texttt{best\_model.zip}. Następnie właśnie ten najlepszy checkpoint był używany jako wynik końcowy.

\section{Tryb graficzny}
Środowisko obsługuje dwa tryby renderowania:
\begin{itemize}
    \item \texttt{human} -- interaktywne wyświetlanie gry w oknie,
    \item \texttt{rgb\_array} -- generowanie klatek obrazu do dalszego zapisu lub analizy.
\end{itemize}

Do renderowania wykorzystano bibliotekę \texttt{pygame}. Dzięki temu środowisko można uruchomić zarówno w trybie przeznaczonym dla człowieka, jak i w trybie przydatnym przy analizie działania agenta. Spełnia to wymaganie rozszerzenia projektu o warstwę graficzną.

\section{Eksperymenty}
Uczenie przeprowadzono dla limitu \texttt{max\_steps = 500}. Według aktualnego podsumowania treningu agent rozegrał 959 epizodów, a średnia nagroda z ostatnich 20 epizodów wyniosła -0.5925. Najlepszy wynik ewaluacyjny uzyskany w trakcie treningu miał średnią nagrodę 8.439.

Końcowy model został porównany z bazą losową. Aktualna ewaluacja została wykonana dla 100 epizodów dla modelu oraz 100 epizodów dla polityki losowej. Wyniki zestawiono w tabeli~\ref{tab:wyniki}.

\begin{table}[H]
\centering
\caption{Porównanie wyuczonego agenta z bazą losową}
\label{tab:wyniki}
\begin{tabular}{lcc}
\toprule
Metryka & Agent DQN & Polityka losowa \\
\midrule
Win rate & 0.49 & 0.00 \\
Collision rate & 0.50 & 0.77 \\
Truncation rate & 0.01 & 0.23 \\
Średnia nagroda & 4.9868 & --- \\
Średnia liczba kroków / czas przeżycia & 103.28 & 311.53 \\
Udział akcji \texttt{up} & 0.487 & 0.204 \\
Udział akcji \texttt{wait} & 0.341 & 0.199 \\
Wait przy ryzykownym ruchu naprzód & 0.579 & 0.201 \\
\bottomrule
\end{tabular}
\end{table}

Najważniejszy wniosek jest taki, że agent uczony metodą DQN nie działa losowo, ponieważ w 49\% epizodów dociera do celu, podczas gdy polityka losowa nie wygrała ani razu. Dodatkowo współczynnik kolizji modelu jest wyraźnie niższy niż dla bazy losowej, a średnia nagroda modelu jest dodatnia. Oznacza to, że agent nauczył się częściowo rozwiązywać zadanie i robi to zauważalnie lepiej niż polityka przypadkowa.

Istotna jest także metryka \textit{wait when forward risky rate}: model czeka w około 57.9\% sytuacji, w których ruch do przodu był ryzykowny, natomiast polityka losowa robi to tylko w około 20.1\% takich przypadków. Rozkład akcji także wskazuje na zachowanie nielosowe. Agent wyraźnie preferuje ruch w górę, ale często korzysta też z oczekiwania oraz sporadycznych ruchów bocznych i w dół. W porównaniu z polityką losową nie rozkłada akcji równomiernie, tylko wykazuje preferencje odpowiadające próbie osiągnięcia celu.

Warto jednak ostrożnie interpretować średni czas przeżycia. Polityka losowa ma tu wyższą wartość głównie dlatego, że w wielu epizodach nie kończy gry kolizją ani sukcesem, tylko dochodzi do limitu 500 kroków i zostaje ucięta. Z tego powodu dla tego zadania bardziej miarodajny jest \textit{win rate} oraz analiza decyzji podejmowanych w sytuacjach ryzyka niż sama liczba kroków.

\section{Wnioski końcowe}
Projekt zakończył się powodzeniem. Udało się przygotować własne środowisko Gymnasium o bardziej złożonej strukturze niż prosta plansza dyskretna, wytrenować agenta oraz dodać tryb graficzny. Otrzymane wyniki pokazują, że agent nauczył się użytecznej, nielosowej polityki, potrafi ukończyć niemal połowę epizodów i osiąga wyraźnie lepsze wyniki niż baza losowa.

W przyszłości środowisko można rozwijać dalej, na przykład przez zwiększenie liczby typów przeszkód, dodanie dynamicznie zmieniających się pasów lub porównanie DQN z innymi algorytmami uczenia ze wzmocnieniem, takimi jak PPO. Już w obecnej postaci projekt spełnia jednak wymagania zadania i pokazuje pełny proces: od zaprojektowania środowiska, przez uczenie agenta, aż po ocenę jakości wyuczonej polityki.

\end{document}
